\documentclass[10pt,a4paper]{amsart}
\usepackage{amsmath,amssymb,amsthm,graphicx}
% Retained preamble for the Stacks Project chapter {CH} (shared project preamble).
% The chapter file pinned-original-source.tex only carries \input{{preamble}} on its first line,
% so no preamble text lives inside the chapter itself.  The Stacks Project uses one single shared
% preamble for every chapter: the snapshot stages it three times (sources/stacks-chow/preamble.tex,
% sources/stacks-etale-cohomology/preamble.tex and focused-seeds/AGP/sources/stacks-snapshot-20261011/
% preamble.tex) and those copies are identical in size.  The block below is copied mechanically,
% line for line, from sources/stacks-chow/preamble.tex lines 14-18 and 154-217: the Xy-pic setup and
% the native theorem declarations together with the mathematical macros.  Nothing was added, removed
% or redefined; the wrapper class and the amsmath/amssymb/amsthm/graphicx base packages are supplied
% by the assembler.
% For commutative diagrams we use Xy-pic
% Chapter: stacks-spaces-more-morphisms.tex -- More on Morphisms of Spaces.
\usepackage[all]{xy}

% We use 2cell for 2-commutative diagrams.
\xyoption{2cell}

% Theorem environments.
%
\theoremstyle{plain}
\newtheorem{theorem}[subsection]{Theorem}
\newtheorem{proposition}[subsection]{Proposition}
\newtheorem{lemma}[subsection]{Lemma}

\theoremstyle{definition}
\newtheorem{definition}[subsection]{Definition}
\newtheorem{example}[subsection]{Example}
\newtheorem{exercise}[subsection]{Exercise}
\newtheorem{situation}[subsection]{Situation}

\theoremstyle{remark}
\newtheorem{remark}[subsection]{Remark}
\newtheorem{remarks}[subsection]{Remarks}

\numberwithin{equation}{subsection}

% Macros
%
\def\lim{\mathop{\mathrm{lim}}\nolimits}
\def\colim{\mathop{\mathrm{colim}}\nolimits}
\def\Spec{\mathop{\mathrm{Spec}}}
\def\Hom{\mathop{\mathrm{Hom}}\nolimits}
\def\Ext{\mathop{\mathrm{Ext}}\nolimits}
\def\SheafHom{\mathop{\mathcal{H}\!\mathit{om}}\nolimits}
\def\SheafExt{\mathop{\mathcal{E}\!\mathit{xt}}\nolimits}
\def\Sch{\mathit{Sch}}
\def\Mor{\mathop{\mathrm{Mor}}\nolimits}
\def\Ob{\mathop{\mathrm{Ob}}\nolimits}
\def\Sh{\mathop{\mathit{Sh}}\nolimits}
\def\NL{\mathop{N\!L}\nolimits}
\def\CH{\mathop{\mathrm{CH}}\nolimits}
\def\proetale{{pro\text{-}\acute{e}tale}}
\def\etale{{\acute{e}tale}}
\def\QCoh{\mathit{QCoh}}
\def\Ker{\mathop{\mathrm{Ker}}}
\def\Im{\mathop{\mathrm{Im}}}
\def\Coker{\mathop{\mathrm{Coker}}}
\def\Coim{\mathop{\mathrm{Coim}}}

% Boxtimes
%
\DeclareMathSymbol{\boxtimes}{\mathbin}{AMSa}{"02}

%
% Macros for moduli stacks/spaces
%
\def\QCohstack{\mathcal{QC}\!\mathit{oh}}
\def\Cohstack{\mathcal{C}\!\mathit{oh}}
\def\Spacesstack{\mathcal{S}\!\mathit{paces}}
\def\Quotfunctor{\mathrm{Quot}}
\def\Hilbfunctor{\mathrm{Hilb}}
\def\Curvesstack{\mathcal{C}\!\mathit{urves}}
\def\Polarizedstack{\mathcal{P}\!\mathit{olarized}}
\def\Complexesstack{\mathcal{C}\!\mathit{omplexes}}
% \Pic is the operator that assigns to X its picard group, usage \Pic(X)
% \Picardstack_{X/B} denotes the Picard stack of X over B
% \Picardfunctor_{X/B} denotes the Picard functor of X over B
\def\Pic{\mathop{\mathrm{Pic}}\nolimits}
\def\Picardstack{\mathcal{P}\!\mathit{ic}}
\def\Picardfunctor{\mathrm{Pic}}
\def\Deformationcategory{\mathcal{D}\!\mathit{ef}}
\newenvironment{reference}{}{}
\newenvironment{slogan}{}{}
\newenvironment{history}{}{}
\usepackage{cleveref}
\crefname{originallemma}{original lemma}{original lemmas}
\Crefname{originallemma}{Original lemma}{Original lemmas}
\makeatletter
\newlabel{decent-spaces-lemma-generalizations-lift-flat}{{decent-spaces-lemma-generalizations-lift-flat}{source}{}{sourcepointer-1}{}}
\newlabel{decent-spaces-lemma-generalizations-lift-flat@cref}{{[originallemma][0][]decent-spaces-lemma-generalizations-lift-flat}{[1][0][]source}{}{}{}}
\makeatother
\crefname{originallemma}{original lemma}{original lemmas}
\Crefname{originallemma}{Original lemma}{Original lemmas}
\makeatletter
\newlabel{spaces-descent-lemma-descending-property-isomorphism}{{spaces-descent-lemma-descending-property-isomorphism}{source}{}{sourcepointer-2}{}}
\newlabel{spaces-descent-lemma-descending-property-isomorphism@cref}{{[originallemma][0][]spaces-descent-lemma-descending-property-isomorphism}{[1][0][]source}{}{}{}}
\makeatother
\crefname{originallemma}{original lemma}{original lemmas}
\Crefname{originallemma}{Original lemma}{Original lemmas}
\makeatletter
\newlabel{spaces-limits-lemma-enough-local}{{spaces-limits-lemma-enough-local}{source}{}{sourcepointer-3}{}}
\newlabel{spaces-limits-lemma-enough-local@cref}{{[originallemma][0][]spaces-limits-lemma-enough-local}{[1][0][]source}{}{}{}}
\makeatother
\crefname{originallemma}{original lemma}{original lemmas}
\Crefname{originallemma}{Original lemma}{Original lemmas}
\makeatletter
\newlabel{spaces-properties-lemma-subscheme}{{spaces-properties-lemma-subscheme}{source}{}{sourcepointer-4}{}}
\newlabel{spaces-properties-lemma-subscheme@cref}{{[originallemma][0][]spaces-properties-lemma-subscheme}{[1][0][]source}{}{}{}}
\makeatother
\title{A flat monomorphism of algebraic spaces is representable by schemes}
\author{The Stacks Project Authors}
\date{stacks-project:spaces-more-morphisms.tex:local-snapshot-20261011; GNU FDL 1.2 or later}
\begin{document}
\maketitle
\noindent\textit{Editorial note.} Statement span is lines 286-289 of sources/stacks-spaces-more-morphisms/pinned-original-source.tex; proof span is lines 291-313, the complete author proof (21 body lines).\par\medskip
\noindent\textit{Editorial note.} Navigation stubs: 4 cited-but-not-included labels are supplied as non-load-bearing pointers with their pinned tags where the snapshot tag index has them.\par\medskip
\noindent\textit{Editorial note.} extra\_packages is empty on purpose: the retained preamble already loads xy with options and a second declaration clashes in the assembler wrapper.\par\medskip
\noindent\textit{Editorial note.} Text is retained verbatim, including the author's own typos and placeholders.\par\medskip
\noindent\textit{Allowed standard prerequisites (editorial).}\begin{itemize}
\item Explicit prerequisite boundary. Flat monomorphisms of algebraic spaces, representability by schemes, and the local structure of a flat monomorphism; the reduction is retained.
\end{itemize}\medskip
\ifdefined\hypertarget\hypertarget{sourcepointer-1}{}\fi
\noindent\textit{Source reference.} Original lemma decent-spaces-lemma-generalizations-lift-flat: sources/stacks-spaces-more-morphisms/pinned-original-source.tex, source lines 313--313. Cited by the retained statement or proof but not part of the collected spans; supplied only as a navigation pointer. The argument does not depend on its content.\par\medskip
\ifdefined\hypertarget\hypertarget{sourcepointer-2}{}\fi
\noindent\textit{Source reference.} Original lemma spaces-descent-lemma-descending-property-isomorphism: sources/stacks-spaces-more-morphisms/pinned-original-source.tex, source lines 313--313. Cited by the retained statement or proof but not part of the collected spans; supplied only as a navigation pointer. The argument does not depend on its content.\par\medskip
\ifdefined\hypertarget\hypertarget{sourcepointer-3}{}\fi
\noindent\textit{Source reference.} Original lemma spaces-limits-lemma-enough-local: sources/stacks-spaces-more-morphisms/pinned-original-source.tex, source lines 313--313. Cited by the retained statement or proof but not part of the collected spans; supplied only as a navigation pointer. The argument does not depend on its content.\par\medskip
\ifdefined\hypertarget\hypertarget{sourcepointer-4}{}\fi
\noindent\textit{Source reference.} Original lemma spaces-properties-lemma-subscheme: sources/stacks-spaces-more-morphisms/pinned-original-source.tex, source lines 313--313. Cited by the retained statement or proof but not part of the collected spans; supplied only as a navigation pointer. The argument does not depend on its content.\par\medskip
\begin{lemma}[David Rydh]
\label{lemma-flat-case}
A flat monomorphism of algebraic spaces is representable by schemes.
\end{lemma}

\begin{proof}
Let $f : X \to Y$ be a flat monomorphism of algebraic spaces.
To prove $f$ is representable, we have to show
$X \times_Y V$ is a scheme for every scheme $V$ mapping to $Y$.
Since being a scheme is local (Properties of Spaces, 
Lemma \ref{spaces-properties-lemma-subscheme}), we may
assume $V$ is affine. Thus we may assume $Y = \Spec(B)$
is an affine scheme. Next, we can assume that $X$ is quasi-compact
by replacing $X$ by a quasi-compact open. The space $X$ is
separated as $X \to X \times_{\Spec(B)} X$ is an isomorphism.
Applying Limits of Spaces, Lemma \ref{spaces-limits-lemma-enough-local}
we reduce to the case where $B$ is local, $X \to \Spec(B)$ is a
flat monomorphism, and
there exists a point $x \in X$ mapping to the closed point of $\Spec(B)$.
Then $X \to \Spec(B)$ is surjective as generalizations
lift along flat morphisms of separated algebraic spaces, see
Decent Spaces, Lemma \ref{decent-spaces-lemma-generalizations-lift-flat}.
Hence we see that $\{X \to \Spec(B)\}$ is an fpqc cover.
Then $X \to \Spec(B)$ is a morphism which becomes an isomorphism
after base change by $X \to \Spec(B)$. Hence it is an isomorphism by
fpqc descent, see Descent on Spaces, Lemma
\ref{spaces-descent-lemma-descending-property-isomorphism}.
\end{proof}

\end{document}
